\section{Limitations}
\label{sec:limitations}

Our fitting schedule is PHerc0139-only and may encode scroll-specific material
or acquisition properties.
Fine-teacher predictions are model-derived soft supervision rather than human
ground truth.
Two held-out calibration scrolls, 16 locked slices, and six blind cubes provide
guards but do not establish broad robustness.
The M7 trust projection is saturated, and its radius was selected under an
earlier objective rather than the final loss in Eq.~\ref{eq:objective}.
Our overlap-calibration sweep also fails to bracket its optimum; threshold 0.45
is a morphology-gated operating choice, not the maximizer of that curve. Human
review is necessarily less reducible than a scalar and should be repeated by
independent readers before a public model claim.

The line fitter has a different failure mode.
A completion may be locally smooth and persistent across slices but still join
the wrong wrap near the core.
Radial, third-component, and cross-sheet gates reduce this risk but cannot prove
global correctness.
The postprocess must remain independently measured and its additions must remain
recoverable during qualification.

\begin{figure}[t]
  \centering
  \figureasset{figure_4_failure_cases}{0.94\columnwidth}{1.45in}
  \caption{Reserved failure cases: faint-sheet loss, residual blobs,
  cross-wrap risk, and examples whose morphology changes near the selected
  operating threshold.}
  \Description{Reserved examples of model and postprocess failure cases.}
  \label{fig:limitations}
\end{figure}

Finally, exact replay requires large, separately licensed artifacts.
The training code, recipes, hashes, and export scaffolding are public-repository
material, but the approximately 31 GB validation corpus, teacher atlases, M7
weights, and student checkpoints require dedicated artifact hosting and resolved
upstream terms.
